
quasi-compacts et quasi-s\'epar\'es, \`a morphismes de transition affines, l'application canonique}
\[
\varinjlim \underline{\Pic}_{X/S}(Z_\alpha)\longrightarrow
\underline{\Pic}_{X/S}\bigl(\varprojlim Z_\alpha\bigr)
\tag{3.1.1}
\]
\emph{est bijective.}

\emph{Par cons\'equent} (EGA IV 8.14.2), \emph{si le foncteur de Picard est repr\'esent\'e par un \(S\)-sch\'ema \(\Pic_{X/S}\), alors \(\Pic_{X/S}\) est localement de pr\'esentation finie.}

En effet, un \'el\'ement \(\lambda_\alpha\) de \(\underline{\Pic}_{X/S}(Z_\alpha)\) est repr\'esent\'e par un faisceau inversible \(L'_\alpha\) sur un \(X_{Z'_\alpha}\) o\`u \(Z'_\alpha\to Z_\alpha\) est un morphisme f.p.p.f. Posons \(Z=\varprojlim Z_\alpha\). L'image \(\lambda\) de \(\lambda_\alpha\) dans \(\underline{\Pic}_{X/S}(Z)\) est repr\'esent\'ee par le faisceau inversible \(L'\) sur \(X_{Z'}\), image inverse de \(L_\alpha\), o\`u \(Z'=Z'_\alpha\times_{Z_\alpha}Z\).

Si \(\lambda\) est l'\'el\'ement neutre de \(\underline{\Pic}_{X/S}(Z)\), il existe un morphisme f.p.p.f. \(Z''\to Z'\) tel que l'image inverse de \(L'\) sur \(X_{Z''}\) soit \(\simeq{O}_{X_{Z''}}\). En vertu de (EGA IV 8.8.2, 8.10.5 (vi), 11.2.6 (ii), 8.5.2 (i), et 8.5.2.4), il existe un \(\beta\geq\alpha\) et un morphisme f.p.p.f. \(Z''_\beta\to Z'_\beta\) o\`u \(Z'_\beta=Z'_\alpha\times_{Z_\alpha}Z_\beta\) tel que l'image inverse de \(L'_\alpha\). sur \(X_{Z''_\beta}\) soit \(\simeq{O}_{X_{Z''_\beta}}\). L'image de \(\lambda_\alpha\) est donc d\'ej\`a l'\'el\'ement neutre dans \(\underline{\Pic}_{X/S}(Z_\beta)\). Par suite 3.1.1 est injective.

La surjectivit\'e de 3.1.1 se d\'emontre d'une fa\c{c}on analogue.

\textbf{Proposition 3.2.} \emph{Soient \(S\) un sch\'ema noeth\'erien et \(X\) un \(S\)-sch\'ema propre. On suppose que le foncteur de Picard est repr\'esent\'e par un \(S\)-sch\'ema \(\Pic_{X/S}\). Alors}~:
\begin{enumerate}[label=(\roman*)]
\item \emph{\(\Pic_{X/S}\) est localement de type fini sur \(S\).}
